% ===================================================================
% CHAPTER 3: THEORETICAL FOUNDATIONS
% ===================================================================
\chapter{Теоретски основи}
\label{chap:theory}

Ова поглавје ги изложува четирите теоретски столба на кои почива трудот: федерираното
учење како парадигма за обука, машинското заборавање како барање за одржување,
архитектурата SISA како конкретна конструкција што го овозможува егзактното заборавање, и --- што е подеднакво важно за толкувањето на резултатите --- теоријата на оценување при
изразена класна нерамнотежа. Поглавјето завршува со физичките ограничувања на уредите на
работ, кои ја определуваат мерната рамка.

\section{Федерирано учење и агрегацијата FedAvg}

Федерираното учење \cite{mcmahan2017fedavg} ја превртува вообичаената насока на движење на
податоците: наместо податоците да патуваат кон моделот, моделот патува кон податоците.
Централниот сервер испраќа тековни параметри до избрани клиенти, секој клиент изведува
локална оптимизација врз сопственото множество, и назад се враќаат само ажурираните
параметри. Алгоритмот \textbf{FedAvg} ги спојува тие ажурирања со тежинска средна вредност
пропорционална на бројот на локални примери.

Три својства на FedAvg се одлучувачки за овој труд:

\begin{enumerate}
  \item \textbf{Тежинска агрегација по број на примери.} Клиент со повеќе податоци влијае повеќе. Затоа во експериментот сите четири јазли добиваат идентично по $1\,000\,000$ реда --- за да не се меша ефектот на методот со ефектот на различна тежина.
  \item \textbf{Агрегација без состојба.} Серверот не чува историја на придонеси по клиент, туку само тековните параметри. Ова е причината зошто егзактното \emph{глобално} заборавање од страната на клиентот е невозможно (§3.4).
  \item \textbf{Осетливост на нерамномерно распределени податоци.} Кога локалните множества не се независни и идентично распределени (\emph{non-IID}), локалните оптимуми се разидуваат и просечувањето може да го влоши моделот. FedProx \cite{li2020fedprox} го ублажува тоа со проксимален член што ја казнува оддалеченоста од глобалните параметри.
\end{enumerate}

Кога дел од клиентите се злонамерни, обичната средна вредност е ранлива: еден доволно
голем придонес може да ја повлече агрегацијата произволно. Оттука семејството робусни
правила за агрегација --- \textbf{Krum} \cite{blanchard2017krum}, кое избира едно ажурирање што е
најблиску до мнозинството, и \textbf{скратената средна вредност} (\emph{trimmed mean})
\cite{yin2018trimmedmean}, која ги отфрла екстремните вредности по координата пред
просечувањето. Овие правила се предмет на претходната фаза на истражувањето (\texttt{results/course/})
и служат како контекст; тие штитат од \emph{влијанието} на труењето, но \textbf{не го отстрануваат} --- што е токму разликата меѓу робусноста и заборавањето.

\section{Машинско заборавање: егзактно наспроти приближно}

Машинското заборавање \cite{cao2015unlearning, nguyen2022musurvey} го дефинира барањето
влијанието на определено подмножество податоци да биде отстрането од веќе обучен модел.
Мотивацијата е двојна: \textbf{регулаторна} --- членот 17 од Општата регулатива за заштита на
податоците \cite{gdpr2016} воведува право на бришење, кое во контекст на машинско учење не
се исполнува со бришење на суровите записи --- и \textbf{безбедносна}: ако напаѓач вградил
внимателно конструирани податоци во детектор на аномалии, детекторот мора да ги заборави за
да ја врати својата функција.

Литературата разликува два режима:

\begin{itemize}
  \item \textbf{Егзактно заборавање} гарантира дека добиениот модел е распределбено идентичен со модел обучен од нула без отстранетите податоци. Гаранцијата е докажлива, но вообичаено се плаќа со пресметка или со ограничувања на архитектурата на обуката.
  \item \textbf{Приближно заборавање} се задоволува со тоа влијанието да биде намалено под некој праг, обично преку насочена измена на параметрите. Постапката е побрза, но гаранцијата е статистичка, не структурна.
\end{itemize}

Овој труд работи исклучиво во егзактниот режим --- и двата споредувани пристапа даваат
егзактно локално заборавање. Тоа е намерен избор: кога двете постапки водат до \textbf{иста
гаранција}, разликата меѓу нив е чиста цена, а не компромис меѓу цена и сигурност.

\subsection{Емпириска проверка на заборавањето}

Наспроти конструкциската гаранција стои емпириската проверка. Стандардната алатка е
\textbf{пробата за заклучување членство} (\emph{membership inference}) \cite{shokri2017mia}: ако
моделот сè уште „памети“ определен пример, неговата загуба врз тој пример ќе биде
систематски пониска отколку врз невиден пример. Границите на таа проверка се важни и се
дискутирани во Глава 8 --- модел со мал број параметри едноставно не протекува доволно
сигнал за пробата да биде заклучна, па егзактноста мора да почива на конструкцијата.

\section{SISA: шардирање, изолација, сегментирање и агрегација}

SISA \cite{bourtoule2021sisa} е конструкција што го прави егзактното заборавање евтино така
што однапред го \textbf{ограничува досегот на секој примерок}. Четирите составни идеи се:

\begin{itemize}
  \item \textbf{Шардирање (\emph{sharded}).} Множеството за обука се дели на $S$ дисјунктни шардови. За секој шард се обучува посебен \textbf{составен модел} (\emph{constituent}). Примерок што влегол во еден шард не може да влијае врз ниту еден друг составен модел.
  \item \textbf{Изолација (\emph{isolated}).} Меѓу составните модели нема никаков проток на информации во текот на обуката. Тоа е условот што го прави ограничувањето на досегот докажливо.
  \item \textbf{Сегментирање (\emph{sliced}).} Секој шард дополнително се дели на $R$ сегменти, кои се вградуваат постапно. По секој сегмент состојбата на моделот се зачувува во контролна точка.
  \item \textbf{Агрегација (\emph{aggregated}).} При предвидување, излезите на $S$-те составни модели се спојуваат --- во оригиналната формулација со гласање на мнозинството.
\end{itemize}

При барање за заборавање на примерок, се идентификува шардот што го содржи, се врати
наназад до контролната точка од последниот \textbf{чист} сегмент, и се обработуваат повторно
само сегментите од таа точка наваму. Останатите $S-1$ шардови остануваат недопрени. Очекуваното
забрзување е пропорционално на $S$ и на позицијата на компромитираниот сегмент во редоследот.

Оттука произлегува и \textbf{условот под кој SISA има смисла}: предноста е најголема кога
компромитирањето е \emph{просторно и временски ограничено} --- малку шардови, доцни сегменти. Тоа
е токму моделот на закана „извор компромитиран во моментот $\tau$“. Ако отстранувањето
погоди ран сегмент, или ако е распространето низ сите шардови, повторната обработка се
приближува до целосна повторна обука и предноста исчезнува. Оваа зависност е позната
слабост на конструкцијата и се третира како услов на важење, не се премолчува.

Цената што SISA ја плаќа однапред е \textbf{складирањето}: $S \times R$ контролни точки по
рунда, запишани непрекинато во текот на целата обука, без оглед на тоа дали воопшто ќе
дојде барање за заборавање.

\section{Федерирано заборавање и границата на FedAvg}

Пренесувањето на заборавањето во федерираната поставка \cite{liu2024fusurvey,
romandini2024fusurvey} воведува проблем што не постои централизирано: моделот е
\textbf{заеднички}, а постапката за заборавање се изведува кај \textbf{еден} учесник.

Предложените решенија се распоредуваат по тоа каква состојба чуваат. \textbf{FedEraser}
\cite{liu2021federaser} ги запишува историските ажурирања на клиентите на серверот и потоа
ја реконструира траекторијата без придонесот на отстранетиот клиент. \textbf{FedRecover}
\cite{cao2023fedrecover} слично користи историски информации за да го обнови глобалниот
модел по напад со труење. Друга насока директно ја пресликува SISA во федерираниот
простор: клиентите се групираат во дисјунктни кластери, за секој кластер се обучува
независен модел, и при барање за заборавање се преобучува само засегнатиот кластер
\cite{lin2024codedsharding}.

Сите тие пристапи имаат заедничка претпоставка: \textbf{серверот соработува} --- чува историја,
или ја менува шемата на агрегација. Овој труд намерно ја испитува спротивната поставка:
стандарден FedAvg без состојба, со постапка за заборавање изведена исклучиво \textbf{локално, на
самиот уред}. Последицата е точно одредена и мора да се наведе: локалната состојба на
клиентот станува докажливо чиста, но придонесот што веќе е агрегиран во $w^{(t)}$ не може
да биде отстранет од страната на клиентот. Заборавањето е \textbf{егзактно локално, приближно
глобално}. Глобалното повлекување на влијанието се случува постепено, преку натамошните
рунди во кои клиентот придонесува со исчистена состојба.

\section{Оценување при класна нерамнотежа}

Овој дел е теоретската подлога за наодот што ја носи \hyperref[chap:results]{Глава~\ref*{chap:results}}, и затоа е изложен пред
резултатите, а не по нив.

Нека матрицата на конфузија содржи $TP$, $FP$, $TN$ и $FN$. Вообичаените показатели се

\[
\text{одѕив} = \frac{TP}{TP + FN}, \qquad
\text{прецизност} = \frac{TP}{TP + FP}, \qquad
F_1 = \frac{2 \cdot \text{прецизност} \cdot \text{одѕив}}{\text{прецизност} + \text{одѕив}} .
\]

Сите три ги игнорираат \textbf{точно негативните} ($TN$) --- ниту еден од нив не ја вреднува
способноста на моделот да препознае бениген сообраќај. Врз множество каде класата
\emph{напад} учествува со удел $\pi$, тривијалниот класификатор што предвидува „сè е напад“
постигнува

\[
\text{одѕив} = 1{,}000, \qquad \text{прецизност} = \pi, \qquad F_1 = \frac{2\pi}{1 + \pi}.
\]

При $\pi = 0{,}965$ тоа дава $F_1 \approx 0{,}982$ --- вредност што во извештај изгледа
одлично, а одговара на модел што \textbf{не открил ниту еден бениген тек}. Ова не е хипотетички
ризик; тоа е точно она што се случува во еден од двата споредувани крака
(\hyperref[chap:results]{Глава~\ref*{chap:results}}, §7.5).

Показателите што ја затвораат таа дупка ја вклучуваат и втората колона на матрицата:

\[
\text{специфичност} = \frac{TN}{TN + FP}, \qquad
\text{избалансирана точност} = \frac{\text{одѕив} + \text{специфичност}}{2},
\]

\[
\mathrm{MCC} = \frac{TP \cdot TN - FP \cdot FN}
{\sqrt{(TP{+}FP)(TP{+}FN)(TN{+}FP)(TN{+}FN)}} .
\]

Избалансираната точност изнесува $0{,}50$ за случаен избор и за секој тривијален
класификатор, независно од нерамнотежата. Метјуовиот коефициент на корелација
\cite{chicco2020mcc} е висок само кога моделот е добар во \textbf{сите четири} полиња на
матрицата, и изнесува $0$ за предвидување без врска со вистината; Chicco и Jurman
аргументираат дека токму затоа треба да се претпочита пред точноста и $F_1$.

Конечно, \textbf{ROC-AUC} ја мери способноста за рангирање независно од прагот на одлука. Разликата
меѓу висок ROC-AUC и ниска избалансирана точност е дијагностички значајна: таа кажува дека
информацијата постои во излезот на моделот, но \textbf{прагот е погрешно поставен} --- дефект на
калибрација, а не на учење. Треба да се има предвид дека и самиот ROC-AUC може да остави
преоптимистички впечаток врз силно нерамнотежени множества, поради што Saito и Rehmsmeier
\cite{saito2015prc} препорачуваат тој да се чита заедно со кривата прецизност–одѕив.

\section{Физички ограничувања на уредите на работ}

Последниот столб е хардверски. Уред од класата Raspberry Pi 5 без активно ладење е
ограничен по четири меѓусебно поврзани оски:

\begin{itemize}
  \item \textbf{Топлина.} Чипот наметнува термичко придушување околу 85 °C. Под секое одржливо оптоварување температурата се заситува, па \textbf{врвната температура не е показател што разликува}; показателот е колку \textbf{време} системот поминал придушен и колку тактот опаднал.
  \item \textbf{Пресметка.} Придушувањето ја намалува работната фреквенција, што ја продолжува пресметката, што пак создава дополнителна топлина --- повратна спрега што ги казнува долгите задачи несразмерно повеќе од кратките.
  \item \textbf{Складирање.} SD-картичка од класата A1 е ограничена на редот од 500 случајни запишувања во секунда. За постапка што чува многу контролни точки тоа е потенцијално тесно грло --- потенцијално, бидејќи дали навистина ќе биде зависи од големината на моделот, што овој труд го \textbf{испитува}, а не го претпоставува.
  \item \textbf{Напојување.} Потрошувачката се мери во ват-часови преку надворешен мерач поставен сериски на напојувањето. Мерењето мора да ја одземе основната потрошувачка во мирување за да ја изрази \textbf{маргиналната} цена на работното оптоварување.
\end{itemize}

Овие ограничувања се намерно задржани во експерименталната поставка (вентилаторот е
исклучен, SD-картичката не е заменета со побрз носител) затоа што тие \textbf{се предметот на
мерење}, а не пречка што треба да се заобиколи.

